\section{Spinor and Gamma-Matrix Conventions}
\label{sec:spinor-conventions}

This section records the spinor conventions used throughout the monograph.
The convention is the Weinberg-compatible four-dimensional gamma-matrix
basis, adapted to the mostly-plus metric used in this monograph.  The section
fixes the data that determine signs in the Dirac adjoint, current
normalization, chirality, Majorana conjugation, anomaly traces, and
supersymmetry formulae.  All comparisons below keep the Lorentzian metric
fixed as \(\eta=\operatorname{diag}(-,+,+,+)\).

\section{Four-Dimensional Lorentzian Clifford Algebra}

\begin{definition}[Four-dimensional mostly-plus spinor convention]
\label{def:four-dimensional-mostly-plus-spinor-convention}
The Lorentzian metric is
\[
  \eta_{\mu\nu}=\operatorname{diag}(-,+,+,+),
  \qquad
  \eta^{\mu\rho}\eta_{\rho\nu}=\delta^\mu{}_\nu .
\]
Gamma matrices are endomorphisms of a complex four-dimensional vector space
\(S\) and obey
\begin{equation}
  \{\gamma^\mu,\gamma^\nu\}=2\eta^{\mu\nu}\mathbf 1_S .
  \label{eq:appendix-lorentz-clifford}
\end{equation}
Consequently
\[
  (\gamma^0)^2=-\mathbf 1_S,
  \qquad
  (\gamma^i)^2=\mathbf 1_S,\quad i=1,2,3 .
\]
For a covector \(p_\mu\) we write
\[
  \not p=p_\mu\gamma^\mu .
\]
If \(p^2=\eta^{\mu\nu}p_\mu p_\nu\), then
\[
  \not p^{\,2}=p^2\mathbf 1_S .
\]

The explicit basis used for sign checks in the monograph is
\begin{equation}
  \gamma^0
  =
  -\ii
  \begin{pmatrix}
    0&I_2\\ I_2&0
  \end{pmatrix},
  \qquad
  \gamma^j
  =
  -\ii
  \begin{pmatrix}
    0&\sigma_j\\ -\sigma_j&0
  \end{pmatrix},
  \qquad j=1,2,3,
  \label{eq:appendix-weinberg-gamma-basis}
\end{equation}
where
\[
  \sigma_1=\begin{pmatrix}0&1\\1&0\end{pmatrix},
  \quad
  \sigma_2=\begin{pmatrix}0&-\ii\\ \ii&0\end{pmatrix},
  \quad
  \sigma_3=\begin{pmatrix}1&0\\0&-1\end{pmatrix}.
\]
The hermiticity identities are
\[
  (\gamma^0)^\dagger=-\gamma^0,
  \qquad
  (\gamma^i)^\dagger=\gamma^i .
\]
Define
\begin{equation}
  \beta=\ii\gamma^0,
  \qquad
  \bar\psi=\psi^\dagger\beta .
  \label{eq:appendix-dirac-adjoint}
\end{equation}
Then
\[
  \beta^\dagger=\beta,
  \qquad
  \beta^2=\mathbf 1_S,
  \qquad
  (\gamma^\mu)^\dagger\beta=-\beta\gamma^\mu .
\]
Thus the bilinear \(\bar\psi\psi\) is Lorentz invariant and
\(\bar\psi\gamma^\mu\psi\) transforms as a Lorentz vector, with the Lorentz
action of Definition~\ref{def:spin-representation-and-chirality}.  In these
conventions the free Lorentzian Dirac density used in the monograph is
\[
  \mathcal L_{\rm D}=-\bar\psi(\gamma^\mu\partial_\mu+m)\psi .
\]
For a real Abelian charge \(q\), the Hermitian vector current is
\[
  j^\mu=\ii q\,\bar\psi\gamma^\mu\psi .
\]
The explicit factor of \(\ii\) is not optional: it follows from
\eqref{eq:appendix-dirac-adjoint} and the anti-Hermitian nature of
\(\bar\psi\gamma^\mu\psi\) in the mostly-plus convention.
\end{definition}

\begin{proposition}[Basic Clifford, adjoint, and current identities]
\label{prop:basic-clifford-adjoint-current-identities}
The data in Definition~\ref{def:four-dimensional-mostly-plus-spinor-convention}
obey
\[
  \not p^{\,2}=p^2\mathbf 1_S,
  \qquad
  \beta^\dagger=\beta,
  \qquad
  \beta^2=\mathbf 1_S,
  \qquad
  (\gamma^\mu)^\dagger\beta=-\beta\gamma^\mu .
\]
Consequently the pairing \(s,t\mapsto s^\dagger\beta t\) is the invariant
Dirac pairing used throughout the monograph, and
\(\ii\bar\psi\gamma^\mu\psi\) is Hermitian for an even classical spinor
field or for the corresponding normal-ordered quantum current.
\end{proposition}

\begin{proof}
The square of \(\not p=p_\mu\gamma^\mu\) is
\[
  p_\mu p_\nu\gamma^\mu\gamma^\nu
  =
  \frac12 p_\mu p_\nu\{\gamma^\mu,\gamma^\nu\}
  =
  p_\mu p_\nu\eta^{\mu\nu}\mathbf 1_S.
\]
The explicit matrices give
\((\gamma^0)^\dagger=-\gamma^0\) and
\((\gamma^i)^\dagger=\gamma^i\).  Since \(\beta=\ii\gamma^0\) and
\((\gamma^0)^2=-1\), this gives \(\beta^\dagger=\beta\) and
\(\beta^2=1\).  For \(\mu=0\),
\[
  (\gamma^0)^\dagger\beta=-\gamma^0\ii\gamma^0=\ii
  =
  -\ii\gamma^0\gamma^0=-\beta\gamma^0 ,
\]
and for \(i=1,2,3\),
\[
  (\gamma^i)^\dagger\beta
  =
  \gamma^i\ii\gamma^0
  =
  -\ii\gamma^0\gamma^i
  =
  -\beta\gamma^i .
\]
Thus \(\beta\gamma^\mu\) is anti-Hermitian.  With the adjoint convention
\(\bar\psi=\psi^\dagger\beta\), the bilinear
\(\bar\psi\gamma^\mu\psi\) is anti-Hermitian in the even finite-dimensional
model.  Multiplication by \(\ii\) gives the Hermitian current normalization.
The quantum statement follows after imposing the same adjoint convention on
the operator-valued distribution and using normal ordering when products at
the same point are represented in a regulator.
\end{proof}

\section{Spin Representation and Chirality}

\begin{definition}[Spin representation and chirality]
\label{def:spin-representation-and-chirality}
The spinor Lorentz generators are
\begin{equation}
  S^{\mu\nu}
  =
  -\frac{\ii}{4}[\gamma^\mu,\gamma^\nu].
  \label{eq:appendix-spin-generators}
\end{equation}
For \(\Lambda=\exp\omega\in SO^+(1,3)\), with
\(\omega_{\mu\nu}=-\omega_{\nu\mu}\), the corresponding element in the
spinor representation is
\[
  R(\Lambda)
  =
  \exp\!\left(-\frac{\ii}{2}\omega_{\mu\nu}S^{\mu\nu}\right).
\]
The Clifford relation gives
\[
  [S^{\mu\nu},\gamma^\rho]
  =
  -\ii\gamma^\mu\eta^{\nu\rho}
  +\ii\gamma^\nu\eta^{\mu\rho},
\]
and hence
\[
  R(\Lambda)\gamma^\mu R(\Lambda)^{-1}
  =
  \Lambda^\mu{}_\nu\gamma^\nu .
\]
The \(\beta\)-pairing is invariant:
\[
  R(\Lambda)^\dagger\beta R(\Lambda)=\beta .
\]

The chirality matrix is
\begin{equation}
  \gamma_5=-\ii\gamma^0\gamma^1\gamma^2\gamma^3 .
  \label{eq:appendix-gamma-five}
\end{equation}
It obeys
\[
  \gamma_5^2=\mathbf 1_S,
  \qquad
  \{\gamma_5,\gamma^\mu\}=0,
  \qquad
  [\gamma_5,S^{\mu\nu}]=0 .
\]
In the explicit basis \eqref{eq:appendix-weinberg-gamma-basis},
\[
  \gamma_5=
  \begin{pmatrix}
    I_2&0\\0&-I_2
  \end{pmatrix},
  \qquad
  P_\pm=\frac{1\pm\gamma_5}{2}.
\]
Thus a spinor \(s\in S\) decomposes as
\[
  s=
  \begin{pmatrix}\chi_a\\ \widetilde\chi_{\dot a}\end{pmatrix},
  \qquad
  \chi\in S_+:=P_+S,
  \qquad
  \widetilde\chi\in S_-:=P_-S,
  \qquad
  a,\dot a=1,2 .
\]
The matrix \(\gamma^\mu\) maps \(S_+\) to \(S_-\) and \(S_-\) to \(S_+\).

With the orientation convention
\[
  \epsilon^{0123}=+1,
\]
the four-dimensional trace convention is
\begin{equation}
  \operatorname{tr}
  \left(\gamma_5\gamma^\mu\gamma^\nu\gamma^\rho\gamma^\sigma\right)
  =
  4\ii\,\epsilon^{\mu\nu\rho\sigma}.
  \label{eq:appendix-gamma-five-trace}
\end{equation}
This is the convention used in the anomaly chapter.
\end{definition}

\begin{proposition}[Spin covariance and chirality trace]
\label{prop:spin-covariance-and-chirality-trace}
The generators in Definition~\ref{def:spin-representation-and-chirality}
satisfy
\[
  [S^{\mu\nu},\gamma^\rho]
  =
  -\ii\gamma^\mu\eta^{\nu\rho}
  +\ii\gamma^\nu\eta^{\mu\rho}.
\]
Consequently
\[
  R(\Lambda)\gamma^\mu R(\Lambda)^{-1}
  =
  \Lambda^\mu{}_\nu\gamma^\nu,
  \qquad
  R(\Lambda)^\dagger\beta R(\Lambda)=\beta .
\]
The chirality matrix obeys
\[
  \gamma_5^2=1,\qquad
  \{\gamma_5,\gamma^\mu\}=0,\qquad
  [\gamma_5,S^{\mu\nu}]=0,
\]
and, with \(\epsilon^{0123}=+1\),
\[
  \operatorname{tr}
  (\gamma_5\gamma^\mu\gamma^\nu\gamma^\rho\gamma^\sigma)
  =
  4\ii\,\epsilon^{\mu\nu\rho\sigma}.
\]
\end{proposition}

\begin{proof}
For any Clifford generators,
\[
  [[\gamma^\mu,\gamma^\nu],\gamma^\rho]
  =
  4(\eta^{\nu\rho}\gamma^\mu-\eta^{\mu\rho}\gamma^\nu).
\]
Multiplying by \(-\ii/4\) gives the displayed commutator for
\(S^{\mu\nu}\).  Exponentiating this infinitesimal identity gives the
Clifford covariance of \(R(\Lambda)\gamma^\mu R(\Lambda)^{-1}\).
The infinitesimal derivative of \(R^\dagger\beta R\) is zero because
\((S^{\mu\nu})^\dagger\beta=\beta S^{\mu\nu}\), which follows from
Proposition~\ref{prop:basic-clifford-adjoint-current-identities}; the value
at the identity is \(\beta\).

Let \(\Gamma=\gamma^0\gamma^1\gamma^2\gamma^3\).  Moving each
\(\gamma^\mu\) through the other three Clifford generators shows that
\(\Gamma\) anticommutes with every \(\gamma^\mu\).  Since
\[
  \Gamma^2=(-1)^6(\gamma^0)^2(\gamma^1)^2(\gamma^2)^2(\gamma^3)^2=-1,
\]
\(\gamma_5=-\ii\Gamma\) squares to \(1\).  Anticommutation with every
\(\gamma^\mu\) implies commutation with every even commutator
\([\gamma^\mu,\gamma^\nu]\), hence with \(S^{\mu\nu}\).

The trace with \(\gamma_5\) is totally antisymmetric in the four vector
indices: moving adjacent Clifford generators past one another changes the
sign, while the metric term produced by the anticommutator has a repeated
index and gives zero in the fully antisymmetric part.  It is therefore enough
to evaluate the ordered component:
\[
  \operatorname{tr}(\gamma_5\gamma^0\gamma^1\gamma^2\gamma^3)
  =
  -\ii\operatorname{tr}(\Gamma^2)
  =
  4\ii .
\]
This gives the stated \(\epsilon^{0123}=+1\) normalization.
\end{proof}

\section{Two-Component Blocks and Index Conventions}

\begin{definition}[Two-component block and epsilon conventions]
\label{def:two-component-block-epsilon-conventions}
Let
\[
  \gamma^\mu
  =
  \begin{pmatrix}
    0&\rho^\mu_{+-}\\
    \rho^\mu_{-+}&0
  \end{pmatrix}
\]
with respect to \(S=S_+\oplus S_-\).  In the basis
\eqref{eq:appendix-weinberg-gamma-basis},
\[
  \rho^0_{+-}=-\ii I_2,
  \qquad
  \rho^j_{+-}=-\ii\sigma_j,
  \qquad
  \rho^0_{-+}=-\ii I_2,
  \qquad
  \rho^j_{-+}=+\ii\sigma_j .
\]
Equivalently,
\[
  (\gamma^\mu s)_a=(\rho^\mu_{+-})_a{}^{\dot b}
  \widetilde\chi_{\dot b},
  \qquad
  (\gamma^\mu s)_{\dot a}=(\rho^\mu_{-+})_{\dot a}{}^b\chi_b .
\]

Spinor \(SU(2)\) indices are raised and lowered by
\[
  \chi^a=\epsilon^{ab}\chi_b,
  \qquad
  \widetilde\chi^{\dot a}
  =
  \epsilon^{\dot a\dot b}\widetilde\chi_{\dot b},
  \qquad
  \epsilon^{12}=\epsilon_{12}=1,
  \qquad
  \epsilon^{\dot1\dot2}=\epsilon_{\dot1\dot2}=1 .
\]
With this convention, the two chiral blocks obey the sign relation
\begin{equation}
  (\rho^\mu_{+-})_a{}^{\dot b}
  =
  -(\rho^\mu_{-+})^{\dot b}{}_a
  =
  -\epsilon^{\dot b\dot c}
  (\rho^\mu_{-+})_{\dot c}{}^d
  \epsilon_{da}.
  \label{eq:appendix-rho-sign-relation}
\end{equation}
In lowered form,
\[
  (\rho_\mu)_{a\dot b}(\rho_\nu)^{a\dot b}
  =
  2\eta_{\mu\nu}.
\]
Complex conjugation exchanges the chiral and antichiral components:
\[
  \overline\chi{}^{\dot a}=(\chi_a)^*,
  \qquad
  \overline{\widetilde\chi}{}^{a}
  =
  (\widetilde\chi_{\dot a})^* .
\]
\end{definition}

\begin{proposition}[Two-component sign and contraction identities]
\label{prop:two-component-sign-contraction-identities}
With the epsilon conventions of
Definition~\ref{def:two-component-block-epsilon-conventions},
\[
  (\rho^\mu_{+-})_a{}^{\dot b}
  =
  -\epsilon^{\dot b\dot c}
  (\rho^\mu_{-+})_{\dot c}{}^d
  \epsilon_{da},
\]
and
\[
  (\rho_\mu)_{a\dot b}(\rho_\nu)^{a\dot b}
  =
  2\eta_{\mu\nu}.
\]
\end{proposition}

\begin{proof}
Let
\[
  \epsilon=
  \begin{pmatrix}0&1\\-1&0\end{pmatrix}.
\]
The matrices in the two chiral blocks are
\[
  \rho^0_{+-}=-\ii I_2,\quad \rho^j_{+-}=-\ii\sigma_j,
  \qquad
  \rho^0_{-+}=-\ii I_2,\quad \rho^j_{-+}=+\ii\sigma_j .
\]
Substituting the Pauli matrices gives the finite matrix identity
\[
  \rho^\mu_{+-}
  =
  \left(-\epsilon\,\rho^\mu_{-+}\,\epsilon\right)^T,
\]
which is precisely the displayed index formula.  For the contraction, lower
the vector index with \(\eta_{\mu\nu}\) and raise the two spinor indices with
the same epsilon matrix:
\[
  M^{a\dot b}=\epsilon^{ac}\epsilon^{\dot b\dot d}M_{c\dot d}.
\]
Direct multiplication of the four displayed \(2\times2\) matrices gives
\[
  \sum_{a,\dot b}(\rho_\mu)_{a\dot b}(\rho_\nu)^{a\dot b}
  =
  -2\,\delta_{\mu0}\delta_{\nu0}
  +2\sum_{j=1}^3\delta_{\mu j}\delta_{\nu j}
  =
  2\eta_{\mu\nu}.
\]
\end{proof}

\section{Majorana Conjugation}

\begin{definition}[Majorana conjugation in the Weinberg-compatible basis]
\label{def:majorana-conjugation-weinberg-basis}
A Majorana condition is a Lorentz-covariant reality condition on the complex
spinor representation.  In the explicit basis
\eqref{eq:appendix-weinberg-gamma-basis}, set
\begin{equation}
  B=\gamma_2 .
  \label{eq:appendix-majorana-B}
\end{equation}
Here \(\gamma_2=\eta_{2\nu}\gamma^\nu=\gamma^2\).  Directly from the Pauli
matrices,
\[
  (\gamma^\mu)^*=B\gamma^\mu B^{-1},
  \qquad
  B^*B=\mathbf 1_S .
\]
For a spinor \(s\), define
\[
  s^c=B^{-1}s^* .
\]
Then \(s^c\) transforms in the same Lorentz representation as \(s\), and
\((s^c)^c=s\).  A Majorana spinor is a spinor satisfying
\[
  s=s^c,
  \qquad\hbox{equivalently}\qquad
  s^*=Bs .
\]
If \(\chi=P_+s\) is the chiral component, then the Majorana spinor may be
written
\[
  s=\chi+B^*\chi^* .
\]
The charge-conjugation matrix in the convention often used for transposed
spinor contractions is
\[
  C=B\gamma^0 .
\]
For a Majorana spinor \(s\),
\[
  \bar s=s^\dagger\ii\gamma^0
  =
  \ii s^T C .
\]
For two Grassmann-even Majorana spinors \(s_1,s_2\), with chiral components
\(\chi_1,\chi_2\), the Lorentz vector bilinear decomposes as
\[
  \bar s_1\gamma^\mu s_2
  =
  \overline\chi_1\,\rho^\mu_{+-}\chi_2
  +
  \overline\chi_2\,\rho^\mu_{+-}\chi_1,
\]
where the precise matrix placement is the one fixed in
\eqref{eq:appendix-rho-sign-relation}.  For Grassmann-odd spinors the same
formula must be combined with the sign rules for moving odd coefficients
through one another.
\end{definition}

\begin{proposition}[Majorana conjugation is Lorentz covariant]
\label{prop:majorana-conjugation-lorentz-covariant}
The anti-linear map \(s\mapsto s^c=B^{-1}s^*\) obeys
\[
  (s^c)^c=s,
  \qquad
  (R(\Lambda)s)^c=R(\Lambda)s^c .
\]
The Majorana condition \(s=s^c\) is therefore a Lorentz-covariant real
structure on the complex spinor representation.  If \(s\) is Majorana, then
\[
  \bar s=\ii s^T C,\qquad C=B\gamma^0 .
\]
\end{proposition}

\begin{proof}
The explicit matrices give
\[
  (\gamma^\mu)^*=B\gamma^\mu B^{-1},
  \qquad
  B^*B=1.
\]
The second identity gives
\[
  (s^c)^c
  =
  B^{-1}(B^{-1}s^*)^*
  =
  B^{-1}(B^*)^{-1}s
  =
  s.
\]
The first identity implies
\((S^{\mu\nu})^*=-B S^{\mu\nu}B^{-1}\).  The minus sign is exactly canceled
by complex conjugating the factor \(-\ii\) in the exponential, hence
\(R(\Lambda)^*=B R(\Lambda)B^{-1}\).  Therefore
\[
  (R(\Lambda)s)^c
  =
  B^{-1}R(\Lambda)^*s^*
  =
  R(\Lambda)B^{-1}s^*
  =
  R(\Lambda)s^c .
\]
If \(s=s^c\), then \(s^*=Bs\), and
\[
  \bar s=s^\dagger\ii\gamma^0
  =
  (s^*)^T\ii\gamma^0
  =
  s^T B^T\ii\gamma^0.
\]
For the displayed basis \(B^T=B\), so the last expression is
\(\ii s^T B\gamma^0=\ii s^T C\).
\end{proof}

\section{Comparison with Wess-Bagger-Type Conventions}

At fixed Lorentzian metric \(\eta=\operatorname{diag}(-,+,+,+)\), a common
chiral gamma-matrix basis, used for example in Wess-Bagger-style supersymmetry
formulae, is
\begin{equation}
  \gamma^0_{\rm WB}
  =
  \begin{pmatrix}
    0&I_2\\ -I_2&0
  \end{pmatrix},
  \qquad
  \gamma^j_{\rm WB}
  =
  \begin{pmatrix}
    0&\sigma_j\\ \sigma_j&0
  \end{pmatrix}.
  \label{eq:appendix-wb-gamma-basis}
\end{equation}
These matrices obey the same Clifford relation
\[
  \{\gamma^\mu_{\rm WB},\gamma^\nu_{\rm WB}\}=2\eta^{\mu\nu}.
\]
The Weinberg-compatible basis used in the monograph is related to
\eqref{eq:appendix-wb-gamma-basis} by the chiral phase change
\begin{equation}
  U_{\rm ch}
  =
  \begin{pmatrix}
    I_2&0\\0&\ii I_2
  \end{pmatrix},
  \qquad
  \gamma^\mu=U_{\rm ch}\gamma^\mu_{\rm WB}U_{\rm ch}^{-1}.
  \label{eq:appendix-wb-phase-change}
\end{equation}
Thus if \(s_{\rm WB}\) and \(s\) are coordinate columns for the same abstract
spinor in the two bases, then
\[
  s=U_{\rm ch}s_{\rm WB}.
\]
The chiral component is unchanged, while the antichiral component is
multiplied by \(\ii\).  This basis change is the source of many sign
differences in two-component supersymmetry formulae.  In particular, with the
index-raising convention used in this section, the monograph has the minus
sign in \eqref{eq:appendix-rho-sign-relation}; in the
Wess-Bagger-type basis one often writes the corresponding relation without
that minus sign, together with a different convention for upper and lower
\(\epsilon\)-symbols and dotted-index contractions.  Formulae copied between
the two systems must therefore translate both the gamma matrices and the
index algebra.

The chirality matrix defined by
\[
  \gamma_5=-\ii\gamma^0\gamma^1\gamma^2\gamma^3
\]
is unchanged by the similarity transformation
\eqref{eq:appendix-wb-phase-change}.  The Dirac adjoint is also the same
abstract Hermitian pairing after the basis change:
\[
  \bar s=s^\dagger(\ii\gamma^0)
  =
  s_{\rm WB}^\dagger
  \left(U_{\rm ch}^\dagger \ii\gamma^0 U_{\rm ch}\right)
  s_{\rm WB}.
\]
The displayed matrix in the last expression is the Wess-Bagger-basis matrix
for the same invariant pairing.

\begin{proposition}[Same-metric Wess-Bagger phase translation]
\label{prop:same-metric-wess-bagger-phase-translation}
At the fixed metric \(\eta=\operatorname{diag}(-,+,+,+)\), the matrix
\[
  U_{\rm ch}=\begin{pmatrix}I_2&0\\0&\ii I_2\end{pmatrix}
\]
relates the Wess-Bagger-type basis
\eqref{eq:appendix-wb-gamma-basis} to the Weinberg-compatible basis by
\[
  \gamma^\mu=U_{\rm ch}\gamma^\mu_{\rm WB}U_{\rm ch}^{-1}.
\]
The chirality operator is unchanged as an endomorphism of the abstract
spinor module, while the antichiral coordinate column is multiplied by
\(\ii\).  Hence any formula imported from Wess-Bagger-type notation must
translate both the gamma matrices and the epsilon-index convention.
\end{proposition}

\begin{proof}
Since \(U_{\rm ch}^{-1}=\operatorname{diag}(I_2,-\ii I_2)\), direct block
multiplication gives
\[
  U_{\rm ch}
  \begin{pmatrix}0&I_2\\-I_2&0\end{pmatrix}
  U_{\rm ch}^{-1}
  =
  -\ii\begin{pmatrix}0&I_2\\I_2&0\end{pmatrix}
  =
  \gamma^0
\]
and, for \(j=1,2,3\),
\[
  U_{\rm ch}
  \begin{pmatrix}0&\sigma_j\\ \sigma_j&0\end{pmatrix}
  U_{\rm ch}^{-1}
  =
  -\ii\begin{pmatrix}0&\sigma_j\\-\sigma_j&0\end{pmatrix}
  =
  \gamma^j.
\]
Both bases are block-off-diagonal with the same chiral splitting, so
\(\gamma_5\) remains \(\operatorname{diag}(I_2,-I_2)\).  The coordinate
relation \(s=U_{\rm ch}s_{\rm WB}\) gives
\((\chi,\widetilde\chi)^T=(\chi_{\rm WB},\ii\widetilde\chi_{\rm WB})^T\),
which is the stated antichiral phase.  The epsilon symbols determine how
lowered and raised two-component indices are identified; therefore the
matrix phase alone is not a complete translation of two-component formulae.
\end{proof}

\section{Dimensional Continuation and \texorpdfstring{\(\gamma_5\)}{gamma5}}

\begin{definition}[Dimensional-reduction chiral trace prescription]
\label{def:dimensional-reduction-chiral-trace-prescription}
Dimensional regularization and dimensional reduction are perturbative
schemes; they are not definitions of a non-perturbative fermionic path
integral.  Whenever a chiral trace is present, the finite-dimensional
four-dimensional data must be specified before continuing loop momenta.
The convention used in the anomaly chapter is:
\[
  \gamma_5=-\ii\gamma^0\gamma^1\gamma^2\gamma^3
\]
is a strictly four-dimensional matrix.  In a split
\(D=4-\varepsilon\) momentum algebra, write
\[
  \ell=\ell_\parallel+\ell_\perp,
\]
where \(\ell_\parallel\) lies in the physical four-dimensional subspace and
\(\ell_\perp\) lies in the continued complement.  Then
\[
  \{\gamma_5,\not\ell_\parallel\}=0,
  \qquad
  [\gamma_5,\not\ell_\perp]=0 .
\]
Together with \eqref{eq:appendix-gamma-five-trace}, these rules define the
chiral trace algebra used in the dimensional-reduction anomaly calculation.
Changing this prescription changes intermediate contact terms and must be
compensated by an explicit finite local counterterm if the same renormalized
Ward identity is to be recovered.
\end{definition}

\section{Three-Dimensional Lorentzian Spinors}

\begin{definition}[Three-dimensional mostly-plus Majorana convention]
\label{def:three-dimensional-mostly-plus-majorana-convention}
For later use in lower-dimensional QFT, take
\[
  \eta_{\mu\nu}=\operatorname{diag}(-,+,+),
  \qquad
  \Gamma^0=\ii\sigma^2,
  \qquad
  \Gamma^1=\sigma^1,
  \qquad
  \Gamma^2=\sigma^3 .
\]
Then
\[
  \{\Gamma^\mu,\Gamma^\nu\}=2\eta^{\mu\nu},
\]
and all three matrices are real.  Therefore a Majorana spinor may be taken
to have real components.  Raise and lower two-component indices by
\[
  \chi^a=\epsilon^{ab}\chi_b,
  \qquad
  \epsilon^{12}=\epsilon_{12}=1 .
\]
The lowered gamma matrices
\[
  \Gamma^\mu_{ab}:=(\Gamma^\mu)_a{}^c\epsilon_{cb}
\]
are symmetric in \(a,b\).  For two Grassmann-even Majorana spinors
\(\zeta_1,\zeta_2\),
\[
  \bar\zeta_1\Gamma^\mu\zeta_2
  =
  \zeta_1^T\ii\Gamma^0\Gamma^\mu\zeta_2
  =
  \ii\,\zeta_1^a\Gamma^\mu_{ab}\zeta_2^b .
\]
\end{definition}

\begin{proposition}[Three-dimensional Clifford and symmetric block identities]
\label{prop:three-dimensional-clifford-symmetric-block-identities}
The matrices in
Definition~\ref{def:three-dimensional-mostly-plus-majorana-convention} obey
\[
  \{\Gamma^\mu,\Gamma^\nu\}=2\eta^{\mu\nu},
\]
and \(\Gamma^\mu_{ab}=(\Gamma^\mu)_a{}^c\epsilon_{cb}\) is symmetric in
\(a,b\).
\end{proposition}

\begin{proof}
The Pauli identities
\(\sigma_i\sigma_j+\sigma_j\sigma_i=2\delta_{ij}\) give
\((\ii\sigma^2)^2=-1\), \((\sigma^1)^2=(\sigma^3)^2=1\), and zero
anticommutators between distinct displayed matrices.  With
\(\epsilon=\begin{pmatrix}0&1\\-1&0\end{pmatrix}\),
\[
  \Gamma^0\epsilon=-I_2,\qquad
  \Gamma^1\epsilon=\begin{pmatrix}-1&0\\0&1\end{pmatrix},
  \qquad
  \Gamma^2\epsilon=\begin{pmatrix}0&1\\1&0\end{pmatrix},
\]
which are all symmetric.
\end{proof}

\section{Two-Dimensional Lorentzian Spinors}

\begin{definition}[Two-dimensional mostly-plus spinor trace convention]
\label{def:two-dimensional-mostly-plus-spinor-trace-convention}
In two-dimensional Lorentzian calculations the monograph uses
\[
  \eta_{\mu\nu}=\operatorname{diag}(-,+),
  \qquad
  \{\gamma^\mu,\gamma^\nu\}=2\eta^{\mu\nu},
  \qquad
  \epsilon^{01}=+1 .
\]
The chirality matrix is
\[
  \gamma=\gamma^0\gamma^1,
  \qquad
  \gamma^2=1,
  \qquad
  \{\gamma,\gamma^\mu\}=0 .
\]
The trace convention compatible with the anomaly calculation is
\[
  \operatorname{tr}(\gamma^\nu\gamma\,\gamma^\rho)
  =
  2\epsilon^{\rho\nu}.
\]
Thus
\[
  \operatorname{tr}(\gamma^\nu\gamma\,\not p)
  =
  2\epsilon^{\rho\nu}p_\rho .
 \]
As in four dimensions, \(\bar\psi=\ii\psi^\dagger\gamma^0\) in the
mostly-plus convention.
\end{definition}

\begin{proposition}[Two-dimensional chirality trace]
\label{prop:two-dimensional-chirality-trace}
The two-dimensional convention of
Definition~\ref{def:two-dimensional-mostly-plus-spinor-trace-convention}
is realized by
\[
  \gamma^0=\ii\sigma^2,\qquad
  \gamma^1=\sigma^1,\qquad
  \gamma=\gamma^0\gamma^1=\sigma^3.
\]
It satisfies
\[
  \operatorname{tr}(\gamma^\nu\gamma\,\gamma^\rho)
  =
  2\epsilon^{\rho\nu}.
\]
\end{proposition}

\begin{proof}
The Clifford relation follows from
\((\ii\sigma^2)^2=-\mathbf1\), \((\sigma^1)^2=\mathbf1\), and
\(\sigma^2\sigma^1=-\sigma^1\sigma^2\), so
\(\{\gamma^\mu,\gamma^\nu\}=2\eta^{\mu\nu}\) in mostly-plus signature.  Also
\[
  \gamma^0\gamma^1=\ii\sigma^2\sigma^1=\sigma^3 .
\]
With \(\epsilon^{01}=+1\), the two nonzero traces are
\[
  \operatorname{tr}(\gamma^0\gamma\gamma^1)=-2,
  \qquad
  \operatorname{tr}(\gamma^1\gamma\gamma^0)=2,
\]
which equal \(2\epsilon^{10}\) and \(2\epsilon^{01}\), respectively; repeated
indices vanish because \(\gamma^\nu\gamma\gamma^\nu\) is proportional to
\(\gamma\), whose trace is zero.
\end{proof}

\begin{definition}[Two-dimensional chiral-component Majorana convention]
\label{def:two-dimensional-chiral-component-majorana-convention}
For comparison with two-dimensional CFT component conventions, it is useful
to use a Majorana basis adapted to light-cone components.  In flat coordinates
\(x^0=t\), \(x^1=\sigma\), set
\[
  \widehat\Gamma^0=-\ii\sigma^1,\qquad
  \widehat\Gamma^1=\sigma^2,
  \qquad
  \widehat\Gamma=\widehat\Gamma^0\widehat\Gamma^1=\sigma^3 .
\]
Let \(\sigma^\pm=\sigma\pm t\), so that
\[
  \widehat\Gamma^+=\widehat\Gamma^1+\widehat\Gamma^0,
  \qquad
  \widehat\Gamma^-=\widehat\Gamma^1-\widehat\Gamma^0.
\]
Spinor indices \(A,B=\pm\) are raised and lowered by
\[
  \epsilon^{+-}=\epsilon_{+-}=1,\qquad
  v^A=\epsilon^{AB}v_B .
\]
The lowered light-cone gamma matrices are
\[
  (\widehat\Gamma^\pm)_{AB}
  :=
  (\widehat\Gamma^\pm)_A{}^C\epsilon_{CB}.
\]
This basis gives
\[
  (\widehat\Gamma^+)_{++}=2\ii,\qquad
  (\widehat\Gamma^-)_{--}=2\ii,
\]
with all other lowered light-cone components zero.  It is the convention that
reduces the covariant Majorana action
\[
  -\frac{\ii}{4\pi}
  \int d^2\sigma\,\sqrt{-g}\,
  \psi\Gamma^a\nabla_a^{\rm spin}\psi
\]
to
\[
  -\frac1{2\pi}\int d^2\sigma\,
  \bigl(\psi_+\partial_-\psi_+
  +\psi_-\partial_+\psi_-\bigr)
\]
in flat coordinates, with the same spinor-index contraction convention.  Its
role in this monograph is translational: it is the convenient bridge to the
chiral free-fermion fields of two-dimensional CFT.  After Wick rotation,
\(\psi_-\) and \(\psi_+\) become the holomorphic and antiholomorphic spin
fields on a Riemann surface, and the global information is carried by a spin
structure rather than by a preferred global spinor frame.  Most
two-dimensional CFT computations are therefore done directly with these
components, their OPEs, and their Szego kernels.  No ghost system or
integration over moduli space is part of this convention; those are separate
structures that are outside the fixed-surface two-dimensional QFT/CFT
framework being recorded here.
\end{definition}

\begin{proposition}[Compatibility of the two-dimensional spinor bases]
\label{prop:two-dimensional-chiral-component-dirac-compatibility}
The Dirac-anomaly basis of
Proposition~\ref{prop:two-dimensional-chirality-trace} and the
Majorana basis of
Definition~\ref{def:two-dimensional-chiral-component-majorana-convention} are
similar representations of the same mostly-plus Clifford algebra.  Explicitly,
with
\[
  U_2=\begin{pmatrix}-\ii&0\\0&1\end{pmatrix},
\]
one has
\[
  \widehat\Gamma^0=U_2\gamma^0U_2^{-1},
  \qquad
  \widehat\Gamma^1=U_2\gamma^1U_2^{-1},
  \qquad
  \widehat\Gamma=U_2\gamma U_2^{-1}.
\]
Consequently the two bases define the same invariant Clifford algebra and the
same chirality splitting, while the free-fermion light-cone component
normalization and the Dirac-anomaly trace normalization are different
coordinate realizations of that algebra.  The monograph uses the
Dirac-anomaly basis when evaluating spinor traces and uses the
chiral-component Majorana basis when comparing to two-dimensional CFT
component conventions.
\end{proposition}

\begin{proof}
Direct multiplication gives
\[
  U_2(\ii\sigma^2)U_2^{-1}=-\ii\sigma^1,
  \qquad
  U_2\sigma^1U_2^{-1}=\sigma^2.
\]
The chirality relation follows by multiplying the two identities.  Lowering
the spinor index of
\[
  \widehat\Gamma^+=
  \begin{pmatrix}0&-2\ii\\0&0\end{pmatrix},
  \qquad
  \widehat\Gamma^-=
  \begin{pmatrix}0&0\\2\ii&0\end{pmatrix}
\]
with
\(\epsilon=\begin{pmatrix}0&1\\-1&0\end{pmatrix}\) gives
\[
  \widehat\Gamma^+\epsilon=
  \begin{pmatrix}2\ii&0\\0&0\end{pmatrix},
  \qquad
  \widehat\Gamma^-\epsilon=
  \begin{pmatrix}0&0\\0&2\ii\end{pmatrix}.
\]
Thus the light-cone components have the stated chiral-component
normalization.  Since \(U_2\) is invertible, all Clifford, chirality, and
trace identities are transported between the two bases after transforming the
coordinate components of the spinors.
\end{proof}

\section{Euclidean Spinors}

\begin{definition}[Recursive Euclidean Clifford basis]
\label{def:recursive-euclidean-clifford-basis}
Euclidean gamma matrices in even dimension \(d\) are denoted
\(\Gamma^1,\ldots,\Gamma^d\) and obey
\[
  \{\Gamma^i,\Gamma^j\}=2\delta^{ij}.
\]
A concrete recursive construction starts in \(d=2\) with
\[
  \Gamma^1_{(2)}=\sigma_1,
  \qquad
  \Gamma^2_{(2)}=\sigma_2,
\]
and for even \(d>2\) sets
\[
  \Gamma^\mu_{(d)}
  =
  \Gamma^\mu_{(d-2)}\otimes(-\sigma_3),
  \qquad
  \mu=1,\ldots,d-2,
\]
\[
  \Gamma^{d-1}_{(d)}
  =
  I_{2^{d/2-1}}\otimes\sigma_1,
  \qquad
  \Gamma^d_{(d)}
  =
  I_{2^{d/2-1}}\otimes\sigma_2 .
\]
The Euclidean chirality matrix is
\[
  \Gamma_E
  =
  \ii^{-d/2}\Gamma^1\cdots\Gamma^d,
  \qquad
  \Gamma_E^2=1 .
\]
One convenient charge-conjugation matrix is
\[
  C_E=\Gamma^1\Gamma^3\cdots\Gamma^{d-1},
  \qquad d\equiv0\pmod4,
\]
and
\[
  C_E=\Gamma_E\Gamma^1\Gamma^3\cdots\Gamma^{d-1},
  \qquad d\equiv2\pmod4.
\]
These formulae fix a basis; invariant statements should be phrased in terms
of the corresponding real, complex, or quaternionic structure on the
Euclidean spinor module.
\end{definition}

\begin{proposition}[Euclidean recursion preserves the Clifford relations]
\label{prop:euclidean-recursion-preserves-clifford-relations}
The recursive construction in
Definition~\ref{def:recursive-euclidean-clifford-basis} gives
\[
  \{\Gamma^i_{(d)},\Gamma^j_{(d)}\}=2\delta^{ij},
  \qquad
  \Gamma_E^2=1
\]
in every even dimension \(d\).
\end{proposition}

\begin{proof}
The \(d=2\) matrices are the Pauli matrices \(\sigma_1,\sigma_2\), so the
claim holds at the initial step.  Assume it holds in dimension \(d-2\).
For \(i,j\le d-2\), tensoring both matrices with \(-\sigma_3\) preserves the
anticommutator because \((-\sigma_3)^2=1\).  The two new matrices
\(I\otimes\sigma_1\) and \(I\otimes\sigma_2\) anticommute with each other and
square to \(1\).  They also anticommute with the old matrices because
\(\sigma_1\) and \(\sigma_2\) anticommute with \(\sigma_3\).  Hence the
Clifford relation holds in dimension \(d\).

Let \(\Omega_d=\Gamma^1\cdots\Gamma^d\).  In Euclidean signature,
\(\Omega_d^2=(-1)^{d(d-1)/2}\).  For even \(d=2r\), this is
\((-1)^r\).  Since \(\Gamma_E=\ii^{-r}\Omega_d\),
\[
  \Gamma_E^2=\ii^{-2r}(-1)^r=(-1)^{-r}(-1)^r=1.
\]
\end{proof}

\section{Six- and Eight-Dimensional Euclidean Blocks}

\begin{definition}[Six- and eight-dimensional Euclidean spinor blocks]
\label{def:six-eight-dimensional-euclidean-spinor-blocks}
For later supersymmetric and internal-symmetry applications, it is useful to
record block conventions for Euclidean \(SO(6)\) and \(SO(8)\) spinors.

In \(d=8\), the chiral and antichiral spinor representations are both real
eight-dimensional.  In a basis in which the charge-conjugation matrix is the
identity, write the off-diagonal gamma-matrix components as
\[
  (\Gamma^i)^{a\dot b}=\mathcal C^{i a\dot b},
  \qquad
  (\Gamma^i)^{\dot b a}=\mathcal C^{i\dot b a},
  \qquad
  \mathcal C^{i\dot b a}=-\mathcal C^{i a\dot b}.
\]
They obey
\[
  \mathcal C^{i a\dot c}\mathcal C^{j\dot c b}
  +
  \mathcal C^{j a\dot c}\mathcal C^{i\dot c b}
  =
  2\delta^{ij}\delta^{ab},
\]
\[
  \mathcal C^{i\dot a c}\mathcal C^{j c\dot b}
  +
  \mathcal C^{j\dot a c}\mathcal C^{i c\dot b}
  =
  2\delta^{ij}\delta^{\dot a\dot b},
\]
\[
  \mathcal C^{i a\dot c}\mathcal C^{i\dot d b}
  +
  \mathcal C^{i b\dot c}\mathcal C^{i\dot d a}
  =
  2\delta^{ab}\delta^{\dot c\dot d}.
\]
The first two identities are the Clifford relations; the third records the
triality symmetry that exchanges vector and spinor representations.

In \(d=6\), use \(\operatorname{Spin}(6)\simeq SU(4)\).  Chiral and
antichiral spinors are the fundamental and antifundamental representations
of \(SU(4)\).  Let \(I,J,K,L=1,\ldots,4\).  The gamma-matrix blocks are
antisymmetric tensors
\[
  \mathcal C^i_{IJ}=-\mathcal C^i_{JI},
  \qquad
  \mathcal C^{iIJ}
  =
  \frac12\epsilon^{IJKL}\mathcal C^i_{KL}
  =
  \left(\mathcal C^i_{JI}\right)^* .
\]
They satisfy
\[
  \sum_{i=1}^6
  \mathcal C^i_{IJ}\mathcal C^{iKL}
  =
  -2
  \left(\delta_I{}^K\delta_J{}^L-\delta_I{}^L\delta_J{}^K\right).
\]
These identities are the ones used when four-dimensional supersymmetry or
higher-dimensional reductions are translated into \(SU(4)\) index notation.
\end{definition}
